\chapter{Abschlie{\ss}ende Arbeiten}\label{chap:AbschliessendeArbeiten}

Zum Schluss muss der Bericht gedruckt und gebunden werden. Jede studentische Arbeit wird mit einem Vortrag abgeschlossen.

	
\section{Drucken und Binden}
Der Bericht ist zweiseitig auf einem Laserdrucker zu drucken. Nur das Titelblatt muss farbig gedruckt werden.  
Der Bericht kann im Kopierlädle auf dem Campus mit einer Kaltbindung (keine Ringbindung, keine Warmbindung) gebunden werden. 
Berücksichtigen Sie, wie viele Exemplare gedruckt werden sollen (siehe Abschnitt~\ref{sec:Einreichen}).

\section{Das Einreichen}\label{sec:Einreichen}
Beim Einreichen ist Folgendes zu beachten:
\begin{enumerate}
 \item Der Bericht ist als PDF am Vortag der Präsentation bei dem Prüfer einzureichen.
 \item Unterzeichnen Sie die \emph{Eigenständigkeitserklärung} in jedem Bericht.
 \item Ein Bericht ist \emph{persönlich} beim Vortrag bei dem Prüfer einzureichen. Falls der Vortrag virtuell stattfindet, kann dies zu einem späteren Zeitpunkt nachgeholt werden.
 \item Jedem anderen Betreuer ist ein Bericht auszuhändigen.
 \item Man sollte auch für sich persönlich einen Bericht erstellen und aufbewahren.
 \item Alle zur Arbeit gehörenden Dateien sind nach Absprache dem Hauptbetreuer zu übergeben.
\end{enumerate}
Bei ihrem zuständigen Prüfer handelt es sich entweder um Prof.~Dr.~R.~I.~Leine oder Prof.~Dr.~C.~D.~Remy.

\section{Der Vortrag}\label{sec:Vortrag}
Die Folien können in \LaTeX, Powerpoint, CorelDRAW oder von
Hand erstellt werden. Eine Vorlage in Powerpoint ist unter \verb"_Vorlagen_INM" abgelegt. Folgende Punkte sind generell zu beachten:
\begin{itemize}
 \item Im Vorfeld des Vortrages ist mit dem Hauptbetreuer abzuklären wie viel Sprechzeit erlaubt ist und wann der Vortrag stattfindet. Die Sprechzeit darf auf keinen Fall überzogen werden.

 \item Machen Sie nicht zu viele Folien. Der Hauptbetreuer soll die Anzahl und die Qualität der Folien überprüfen.

 \item Die Titelfolie beinhaltet mindestens:
  \begin{itemize}
   \item Titel der Studienarbeit
   \item Name des Studierenden
   \item Typ der Arbeit (z.B.~Masterarbeit)
   \item Datum
   \item Universitäts- und Institutslogo
   \item Die Namen aller Betreuer
  \end{itemize}

  \item Achten Sie darauf, dass Ihre Folien (inklusive Grafiken) durch den Beamer gut angezeigt werden. Bearbeiten Sie gegebenenfalls Ihre Folien.
\end{itemize}
